\documentclass[11pt,a4paper]{article}

% ── Packages ──────────────────────────────────────────────────────────
\usepackage[utf8]{inputenc}
\usepackage[T1]{fontenc}
\usepackage{lmodern}
\usepackage[margin=2.5cm]{geometry}
\usepackage{booktabs}
\usepackage{tabularx}
\usepackage{longtable}
\usepackage{array}
\usepackage{xcolor}
\usepackage{listings}
\usepackage{hyperref}
\usepackage{fancyhdr}
\usepackage{titlesec}
\usepackage{enumitem}
\usepackage{amssymb}
\usepackage{amsmath}
\usepackage{graphicx}
\usepackage{float}
\usepackage{caption}
\usepackage{bytefield}
\bytefieldsetup{endianness=big}
\usepackage{multicol}
\usepackage{tcolorbox}
\usepackage{tikz}

% ── Colours ───────────────────────────────────────────────────────────
\definecolor{ee8blue}{RGB}{0,90,160}
\definecolor{ee8gray}{RGB}{100,100,100}
\definecolor{codebg}{RGB}{248,248,248}
\definecolor{codeborder}{RGB}{200,200,200}
\definecolor{keywordcolor}{RGB}{0,100,170}
\definecolor{commentcolor}{RGB}{100,140,100}
\definecolor{stringcolor}{RGB}{170,60,0}
\definecolor{warnbg}{RGB}{255,250,230}
\definecolor{warnborder}{RGB}{220,180,50}
\definecolor{notebg}{RGB}{235,245,255}
\definecolor{noteborder}{RGB}{70,130,200}

% ── Listings style ────────────────────────────────────────────────────
\lstdefinelanguage{ee8asm}{
  morekeywords={ADD,SUB,AND,OR,XOR,SHL,SHR,MOV,LDI,LT,EQ,JMP,JZ,JNZ,HALT,IN,OUT},
  sensitive=true,
  morecomment=[l]{;},
  morestring=[b]",
}

\lstset{
  language=ee8asm,
  basicstyle=\ttfamily\small,
  keywordstyle=\color{keywordcolor}\bfseries,
  commentstyle=\color{commentcolor}\itshape,
  stringstyle=\color{stringcolor},
  backgroundcolor=\color{codebg},
  frame=single,
  rulecolor=\color{codeborder},
  framesep=5pt,
  xleftmargin=10pt,
  xrightmargin=10pt,
  numbers=none,
  breaklines=true,
  showstringspaces=false,
  tabsize=4,
  columns=fullflexible,
  keepspaces=true,
  literate={→}{$\rightarrow$}1 {≠}{$\neq$}1 {↔}{$\leftrightarrow$}1,
}

% ── tcolorbox styles ─────────────────────────────────────────────────
\newtcolorbox{designnote}{
  colback=notebg,
  colframe=noteborder,
  fonttitle=\bfseries,
  title=Design Note,
  boxrule=0.5pt,
  arc=2pt,
  left=6pt, right=6pt, top=4pt, bottom=4pt,
}

\newtcolorbox{warningbox}{
  colback=warnbg,
  colframe=warnborder,
  fonttitle=\bfseries,
  title=Important,
  boxrule=0.5pt,
  arc=2pt,
  left=6pt, right=6pt, top=4pt, bottom=4pt,
}

% ── Hyperref setup ────────────────────────────────────────────────────
\hypersetup{
  colorlinks=true,
  linkcolor=ee8blue,
  urlcolor=ee8blue,
  citecolor=ee8blue,
  pdftitle={EE8 CPU -- Instruction Set Architecture Specification},
  pdfauthor={EG2300 Course Team},
}

% ── Header / footer ──────────────────────────────────────────────────
\pagestyle{fancy}
\fancyhf{}
\fancyhead[L]{\small\textcolor{ee8gray}{EE8 CPU --- ISA Specification v4.1}}
\fancyhead[R]{\small\textcolor{ee8gray}{EG2300}}
\fancyfoot[C]{\small\textcolor{ee8gray}{\thepage}}
\renewcommand{\headrulewidth}{0.4pt}
\renewcommand{\footrulewidth}{0pt}

% ── Section formatting ────────────────────────────────────────────────
\titleformat{\section}
  {\Large\bfseries\color{ee8blue}}{\thesection}{1em}{}[\vspace{-0.5em}\rule{\textwidth}{0.4pt}]
\titleformat{\subsection}
  {\large\bfseries\color{ee8blue!80!black}}{\thesubsection}{1em}{}
\titleformat{\subsubsection}
  {\normalsize\bfseries}{\thesubsubsection}{1em}{}

% ── Custom commands ───────────────────────────────────────────────────
\newcommand{\reg}[1]{\texttt{R#1}}
\newcommand{\opcode}[1]{\texttt{#1}}
\newcommand{\signal}[1]{\textsc{#1}}
\newcommand{\instr}[1]{\texttt{\textbf{#1}}}

% ══════════════════════════════════════════════════════════════════════
\begin{document}

% ── Title page ────────────────────────────────────────────────────────
\begin{titlepage}
\centering
\vspace*{3cm}

{\Huge\bfseries\color{ee8blue} EE8 CPU}\\[0.8cm]
{\LARGE Instruction Set Architecture Specification}\\[1.5cm]

\rule{0.6\textwidth}{1pt}\\[1.5cm]

{\Large Version 4.1}\\[0.6cm]
{\large Second Year Electrical Engineering}\\[0.3cm]
{\large Engineering Design Module (EG2300)}\\[3cm]

\vfill
{\small\textcolor{ee8gray}{The EE8 is a simple 8-bit CPU designed for educational purposes.\\
Students will implement this CPU in Logisim to control a batch diverter manufacturing station.}}
\end{titlepage}

% ── Table of contents ─────────────────────────────────────────────────
\tableofcontents
\newpage

% ══════════════════════════════════════════════════════════════════════
\section{Overview}

The EE8 is a simple 8-bit CPU designed for educational purposes. It features a 16-bit instruction word, 4 general-purpose registers, and a minimal RISC-style instruction set. Students will implement this CPU in Logisim to control a batch diverter manufacturing station.

\subsection{Key Specifications}

\begin{table}[H]
\centering
\begin{tabular}{@{}ll@{}}
\toprule
\textbf{Parameter} & \textbf{Value} \\
\midrule
Data width           & 8 bits \\
Instruction width    & 16 bits \\
Number of registers  & 4 \\
Addressing           & 4-bit (16 instructions max) \\
Program memory       & 16-word ROM \\
Signed representation & Two's complement \\
\bottomrule
\end{tabular}
\end{table}

% ══════════════════════════════════════════════════════════════════════
\section{Registers}

The EE8 has four 8-bit general-purpose registers, addressed with 2 bits.

\begin{table}[H]
\centering
\begin{tabular}{@{}ccl@{}}
\toprule
\textbf{Code} & \textbf{Name} & \textbf{Description} \\
\midrule
\texttt{00} & \reg{0} & General purpose register \\
\texttt{01} & \reg{1} & General purpose register \\
\texttt{10} & \reg{2} & General purpose register \\
\texttt{11} & \reg{3} & General purpose register (also used for comparison results) \\
\bottomrule
\end{tabular}
\end{table}

All registers are fully general-purpose and can be used for computation, storage, or I/O operations.

\subsection{R3 and Comparison Operations}

\reg{3} serves double duty as the destination for comparison instructions:
\begin{itemize}[nosep]
  \item \instr{LT Rs1 Rs2} sets \reg{3} to \texttt{1} if Rs1 $<$ Rs2, otherwise \texttt{0}
  \item \instr{EQ Rs1 Rs2} sets \reg{3} to \texttt{1} if Rs1 $=$ Rs2, otherwise \texttt{0}
  \item \reg{3} can be read and written like any register, but comparison instructions will overwrite it
\end{itemize}

\begin{designnote}
Unlike traditional CPUs (x86, ARM) which use packed bit flags (N, Z, C, V), EE8 uses a full 8-bit register for comparison results. This design choice follows modern RISC philosophy (RISC-V):

\textbf{Advantages:}
\begin{itemize}[nosep]
  \item \textbf{Simpler hardware} --- No bit-field extraction logic needed
  \item \textbf{Orthogonal design} --- \reg{3} can be used like any other register
  \item \textbf{Easier to implement} --- Students build standard 8-bit register, not special flag register
  \item \textbf{Modern approach} --- RISC-V (2010s) uses explicit comparison results, not implicit flags
\end{itemize}

\textbf{Trade-offs:}
\begin{itemize}[nosep]
  \item Cannot branch on overflow or unsigned comparisons without additional instructions
  \item Acceptable for EE8's educational scope and target applications
\end{itemize}
\end{designnote}

% ══════════════════════════════════════════════════════════════════════
\section{Platform Signals \& Safety Logic}

The batch diverter station has safety signals that are enforced in \textbf{hardware}, external to the CPU's instruction execution. The CPU does \textbf{not} read or poll these signals --- they act on the platform directly.

\subsection{Signal Definitions}

\begin{table}[H]
\centering
\begin{tabularx}{\textwidth}{@{}llX@{}}
\toprule
\textbf{Signal} & \textbf{Type} & \textbf{Description} \\
\midrule
\signal{reboot}  & Input (async) & System reset. Clears all CPU state and restarts execution. \\
\signal{stop}    & Input         & Operator halt/pause request. \\
\signal{fault}   & Input         & Non-emergency fault (jam, guard open, etc.). \\
\signal{sys\_err} & Internal (latched) & Illegal instruction detected. Set by decode logic. \\
\signal{cpu\_halt} & Internal (latched) & Set by the \instr{HALT} instruction. \\
\bottomrule
\end{tabularx}
\end{table}

\subsection{REBOOT}

\signal{reboot} is an asynchronous hardware reset signal.

While \signal{reboot} = 1:
\begin{itemize}[nosep]
  \item PC is forced to 0
  \item All registers are cleared to 0
  \item Output port is cleared to 0 (\signal{motor\_en}=0, \signal{gate\_sel}=0)
  \item \signal{sys\_err} latch is cleared
  \item \signal{cpu\_halt} latch is cleared
\end{itemize}

Execution starts from address 0 when \signal{reboot} returns to 0. \signal{reboot} is the \textbf{only} way to clear \signal{sys\_err} and \signal{cpu\_halt}.

\subsection{STOP}

\signal{stop} is a hardware halt/pause signal (operator-initiated).

While \signal{stop} = 1:
\begin{itemize}[nosep]
  \item CPU clock is gated --- PC is frozen, no instructions execute
  \item All registers and output port are preserved (retain last values)
  \item Execution resumes from the paused point when \signal{stop} returns to 0
\end{itemize}

\subsection{FAULT}

\signal{fault} is a hardware halt/pause signal with alarm (fault condition).

While \signal{fault} = 1:
\begin{itemize}[nosep]
  \item CPU clock is gated --- same freeze behaviour as \signal{stop}
  \item \signal{alarm} output is asserted (see Section~\ref{sec:alarm})
  \item Execution resumes from the paused point when \signal{fault} returns to 0
\end{itemize}

\subsection{SYS\_ERR}

\signal{sys\_err} is a latched internal signal set by the instruction decoder on illegal instruction detection.

\textbf{Triggers:}
\begin{itemize}[nosep]
  \item Execution of undefined opcode \texttt{1110} (when no team extension is implemented)
  \item Execution of undefined sub-encoding (\texttt{Rs2=01} in the MOV/IN/OUT family under Option~B)
\end{itemize}

\textbf{Behaviour when set:}
\begin{itemize}[nosep]
  \item \signal{sys\_err} is \textbf{latched} --- once set, it remains asserted until \signal{reboot}
  \item PC is frozen (no further instructions execute)
  \item \signal{motor\_en} is forced to 0 (see Section~\ref{sec:haltlogic})
  \item \signal{alarm} is asserted (see Section~\ref{sec:alarm})
\end{itemize}

\textbf{Clearing:} Only \signal{reboot} clears the \signal{sys\_err} latch.

\subsection{Platform Halt Logic}\label{sec:haltlogic}

The platform combines all halt sources into a single \signal{halted} signal:
\[
  \text{HALTED} = \text{STOP} \;\mathbin{|}\; \text{FAULT} \;\mathbin{|}\; \text{SYS\_ERR} \;\mathbin{|}\; \text{CPU\_HALT}
\]

When \signal{halted} is asserted, the conveyor motor is unconditionally disabled:
\[
  \text{MOTOR\_EN}_{\text{pin}} = \text{MOTOR\_EN}_{\text{cpu}} \;\mathbin{\&}\; \overline{\text{HALTED}}
\]

This means the CPU can \emph{request} \signal{motor\_en}=1 via the output port, but the platform will override it to 0 whenever any halt condition is active. This is a hardware safety interlock.

\subsection{Alarm Logic}\label{sec:alarm}

The \signal{alarm} output is driven by hardware, not by the CPU's output port:
\[
  \text{ALARM} = \text{FAULT} \;\mathbin{|}\; \text{SYS\_ERR}
\]

\signal{alarm} is \textbf{not} asserted for \signal{stop} (operator-initiated, not a fault) or \signal{cpu\_halt} (intentional program termination). \signal{alarm} is not part of the CPU's output port --- teams do not need to set or clear it in software.

\begin{warningbox}
These signals are \textbf{not} part of the ISA. The CPU does not execute instructions to read \signal{stop}, \signal{fault}, or \signal{reboot}. They are external hardware signals that act on the platform and clock/reset logic. Students implement these as combinational logic outside the CPU datapath.
\end{warningbox}

% ══════════════════════════════════════════════════════════════════════
\section{I/O Interface}

The EE8 interfaces with the batch diverter station through input and output ports. Teams must choose \textbf{one} of the following I/O strategies and document their choice in the report.

\subsection{Port Definitions}

\textbf{Input Port:}

\begin{table}[H]
\centering
\begin{tabularx}{\textwidth}{@{}clX@{}}
\toprule
\textbf{Bit} & \textbf{Signal} & \textbf{Description} \\
\midrule
0 & \signal{item\_pulse} & Item detection flag (see Section~\ref{sec:itempulse}) \\
\bottomrule
\end{tabularx}
\end{table}

Bits 1--7 are reserved (read as 0).

\textbf{Output Port:}

\begin{table}[H]
\centering
\begin{tabularx}{\textwidth}{@{}clX@{}}
\toprule
\textbf{Bit} & \textbf{Signal} & \textbf{Description} \\
\midrule
0   & \signal{gate\_sel} & Active packing lane (0 = Gate A, 1 = Gate B) \\
1   & \signal{motor\_en} & Conveyor motor enable request (see Section~\ref{sec:haltlogic}) \\
2+  & \signal{status}    & Debug/status bits (optional, team-defined) \\
\bottomrule
\end{tabularx}
\end{table}

\textbf{Note:} \signal{motor\_en} is a \emph{request} --- the platform overrides it to 0 when \signal{halted} (see Section~\ref{sec:haltlogic}). \signal{alarm} is driven by hardware safety logic (see Section~\ref{sec:alarm}) and is not part of the CPU output port.

\subsection{ITEM\_PULSE Semantics}\label{sec:itempulse}

\signal{item\_pulse} is a \textbf{hardware-latched flag} with the following behaviour:
\begin{itemize}[nosep]
  \item \textbf{Set} by external hardware when an item passes the counting sensor
  \item \textbf{Cleared on read} --- when the CPU reads the input port, the \signal{item\_pulse} latch is automatically reset to 0
\end{itemize}

This ``clear-on-read'' mechanism guarantees exactly one count per item. The CPU does not need to wait for the pulse to go low; each read either sees a 1 (item arrived since last read) or 0 (no item). This prevents multi-counting from polling too fast.

\subsection{Reset State}

On reset (\signal{reboot}), the output port is cleared to 0:
\begin{itemize}[nosep]
  \item \signal{gate\_sel} = 0 (Gate A selected)
  \item \signal{motor\_en} = 0 (motor off)
  \item \signal{status} = 0
\end{itemize}

\subsection{Option A: Memory-Mapped I/O}

Specific addresses in the ROM address space are aliased to I/O ports instead of instructions. This approach requires additional decode logic but no new instructions.

\textbf{Implementation:}
\begin{itemize}[nosep]
  \item Address \texttt{0xE} (14): Reading fetches input port value instead of instruction
  \item Address \texttt{0xF} (15): Writing stores to output port instead of executing
\end{itemize}

\subsection{Option B: Dedicated I/O Instructions}

Add explicit \instr{IN} and \instr{OUT} instructions to the ISA. This is the cleanest approach for educational purposes.

\textbf{New Instructions:}

\begin{table}[H]
\centering
\begin{tabular}{@{}cclll@{}}
\toprule
\textbf{Opcode} & \textbf{Binary} & \textbf{Mnemonic} & \textbf{Format} & \textbf{Operation} \\
\midrule
--- & \texttt{0111} & \instr{IN}  & R-Type variant & Rd = input\_port \\
--- & \texttt{0111} & \instr{OUT} & R-Type variant & output\_port = Rs \\
\bottomrule
\end{tabular}
\end{table}

See Section~\ref{sec:io_instructions} for full instruction definitions.

\subsection{Option C: Register-Mapped I/O}

Dedicate specific registers to I/O. Inputs appear in one register; outputs are driven from another.

\textbf{Implementation:}
\begin{itemize}[nosep]
  \item \textbf{Input Register (read-only view):} Reading \reg{0} returns current input pin state
  \item \textbf{Output Register:} Writing \reg{2} drives output pins
\end{itemize}

% ══════════════════════════════════════════════════════════════════════
\section{Instruction Formats}

All instructions are 16 bits wide. There are three instruction formats.

\subsection{R-Type (Register Operations)}

Used for ALU operations, comparisons, and register moves.

\begin{center}
\begin{bytefield}[bitwidth=1.6em]{16}
  \bitheader{0-15} \\
  \bitbox{4}{\footnotesize opcode} &
  \bitbox{2}{\footnotesize Rs1} &
  \bitbox{2}{\footnotesize Rs2} &
  \bitbox{2}{\footnotesize Rd} &
  \bitbox{6}{\footnotesize unused}
\end{bytefield}
\end{center}

\textbf{Fields:}
\begin{itemize}[nosep]
  \item \texttt{opcode} [15:12] --- Operation to perform
  \item \texttt{Rs1} [11:10] --- First source register
  \item \texttt{Rs2} [9:8] --- Second source register
  \item \texttt{Rd} [7:6] --- Destination register
  \item \texttt{unused} [5:0] --- Reserved (should be set to 0)
\end{itemize}

\textbf{Assembly syntax:} \lstinline|MNEMONIC Rs1 Rs2 Rd|

\textbf{Example:} \lstinline|ADD R0 R1 R2| --- Adds \reg{0} and \reg{1}, stores result in \reg{2}.

\textbf{Example 2 --- Adding a constant:}
To add a constant (e.g., 50) to \reg{0}, first load the constant into a register:
\begin{lstlisting}
LDI R1 50       ; Load constant 50 into R1
ADD R0 R1 R0    ; R0 = R0 + 50
\end{lstlisting}

\subsection{I-Type (Immediate)}

Used for loading immediate values into registers.

\begin{center}
\begin{bytefield}[bitwidth=1.6em]{16}
  \bitheader{0-15} \\
  \bitbox{4}{\footnotesize opcode} &
  \bitbox{2}{\footnotesize Rd} &
  \bitbox{2}{\footnotesize unused} &
  \bitbox{8}{\footnotesize immediate}
\end{bytefield}
\end{center}

\textbf{Fields:}
\begin{itemize}[nosep]
  \item \texttt{opcode} [15:12] --- Operation to perform
  \item \texttt{Rd} [11:10] --- Destination register
  \item \texttt{unused} [9:8] --- Reserved (should be set to 0)
  \item \texttt{immediate} [7:0] --- 8-bit constant value (0--255 unsigned, or $-128$ to 127 signed)
\end{itemize}

\textbf{Assembly syntax:} \lstinline|LDI Rd immediate|

\textbf{Example:} \lstinline|LDI R0 42| --- Loads the value 42 into \reg{0}.

\subsection{J-Type (Jump)}

Used for control flow (jumps and conditional branches).

\begin{center}
\begin{bytefield}[bitwidth=1.6em]{16}
  \bitheader{0-15} \\
  \bitbox{4}{\footnotesize opcode} &
  \bitbox{2}{\footnotesize Rcond} &
  \bitbox{6}{\footnotesize unused} &
  \bitbox{4}{\footnotesize address}
\end{bytefield}
\end{center}

\textbf{Fields:}
\begin{itemize}[nosep]
  \item \texttt{opcode} [15:12] --- Jump type
  \item \texttt{Rcond} [11:10] --- Register to test for conditional branches (JZ/JNZ). Ignored for JMP (don't-care, set to \texttt{00}).
  \item \texttt{unused} [9:4] --- Reserved (should be set to 0)
  \item \texttt{address} [3:0] --- Target instruction address (0--15)
\end{itemize}

\textbf{Assembly syntax:}
\begin{itemize}[nosep]
  \item \lstinline|JMP address| --- Unconditional jump
  \item \lstinline|JZ Rcond address| --- Jump if register Rcond $= 0$
  \item \lstinline|JNZ Rcond address| --- Jump if register Rcond $\neq 0$
\end{itemize}

\textbf{Examples:}
\begin{itemize}[nosep]
  \item \lstinline|JMP 5| --- Jump to instruction at address 5
  \item \lstinline|JNZ R3 5| --- If \reg{3} is not zero, jump to address 5
  \item \lstinline|JZ R0 3| --- If \reg{0} is zero, jump to address 3
\end{itemize}

\begin{designnote}
The Rcond field allows conditional branches to test \textbf{any} register, not just \reg{3}. This makes programs shorter and more flexible --- critical when the ROM is limited to 16 instructions. Comparison results still go to \reg{3}, but you can also branch directly on any register being zero/non-zero (e.g., a loop counter in \reg{0}).
\end{designnote}

% ══════════════════════════════════════════════════════════════════════
\section{Instruction Set Reference}

\subsection{Arithmetic Instructions}

\subsubsection{ADD --- Add}

\begin{table}[H]
\begin{tabularx}{\textwidth}{@{}lX@{}}
\toprule
\textbf{Opcode}      & \opcode{0000} \\
\textbf{Format}      & R-Type \\
\textbf{Syntax}      & \lstinline|ADD Rs1 Rs2 Rd| \\
\textbf{Operation}   & \texttt{Rd = Rs1 + Rs2} \\
\textbf{Description} & Adds the values in Rs1 and Rs2, stores the result in Rd. Overflow wraps around (modulo 256). \\
\bottomrule
\end{tabularx}
\end{table}

\begin{lstlisting}
LDI R0 10      ; R0 = 10
LDI R1 20      ; R1 = 20
ADD R0 R1 R2   ; R2 = 30
\end{lstlisting}

\subsubsection{SUB --- Subtract}

\begin{table}[H]
\begin{tabularx}{\textwidth}{@{}lX@{}}
\toprule
\textbf{Opcode}      & \opcode{0001} \\
\textbf{Format}      & R-Type \\
\textbf{Syntax}      & \lstinline|SUB Rs1 Rs2 Rd| \\
\textbf{Operation}   & \texttt{Rd = Rs1 - Rs2} \\
\textbf{Description} & Subtracts Rs2 from Rs1, stores the result in Rd. Uses two's complement for negative results. \\
\bottomrule
\end{tabularx}
\end{table}

\begin{lstlisting}
LDI R0 50      ; R0 = 50
LDI R1 30      ; R1 = 30
SUB R0 R1 R2   ; R2 = 20
\end{lstlisting}

\subsection{Bitwise Instructions}

\subsubsection{AND --- Bitwise AND}

\begin{table}[H]
\begin{tabularx}{\textwidth}{@{}lX@{}}
\toprule
\textbf{Opcode}      & \opcode{0010} \\
\textbf{Format}      & R-Type \\
\textbf{Syntax}      & \lstinline|AND Rs1 Rs2 Rd| \\
\textbf{Operation}   & \texttt{Rd = Rs1 \& Rs2} \\
\textbf{Description} & Performs bitwise AND on Rs1 and Rs2. \\
\bottomrule
\end{tabularx}
\end{table}

\begin{lstlisting}
LDI R0 0b11110000   ; R0 = 240
LDI R1 0b10101010   ; R1 = 170
AND R0 R1 R2        ; R2 = 0b10100000 = 160
\end{lstlisting}

\subsubsection{OR --- Bitwise OR}

\begin{table}[H]
\begin{tabularx}{\textwidth}{@{}lX@{}}
\toprule
\textbf{Opcode}      & \opcode{0011} \\
\textbf{Format}      & R-Type \\
\textbf{Syntax}      & \lstinline|OR Rs1 Rs2 Rd| \\
\textbf{Operation}   & \texttt{Rd = Rs1 | Rs2} \\
\textbf{Description} & Performs bitwise OR on Rs1 and Rs2. \\
\bottomrule
\end{tabularx}
\end{table}

\subsubsection{XOR --- Bitwise Exclusive OR}

\begin{table}[H]
\begin{tabularx}{\textwidth}{@{}lX@{}}
\toprule
\textbf{Opcode}      & \opcode{0100} \\
\textbf{Format}      & R-Type \\
\textbf{Syntax}      & \lstinline|XOR Rs1 Rs2 Rd| \\
\textbf{Operation}   & \texttt{Rd = Rs1 \^{} Rs2} \\
\textbf{Description} & Performs bitwise XOR on Rs1 and Rs2. \\
\bottomrule
\end{tabularx}
\end{table}

\subsection{Shift Instructions}

\subsubsection{SHL --- Shift Left}

\begin{table}[H]
\begin{tabularx}{\textwidth}{@{}lX@{}}
\toprule
\textbf{Opcode}      & \opcode{0101} \\
\textbf{Format}      & R-Type \\
\textbf{Syntax}      & \lstinline|SHL Rs1 Rs2 Rd| \\
\textbf{Operation}   & \texttt{Rd = Rs1 <{}< Rs2} \\
\textbf{Description} & Shifts Rs1 left by Rs2 bits. Vacated bits are filled with zeros. Bits shifted out are discarded. \\
\bottomrule
\end{tabularx}
\end{table}

\begin{lstlisting}
LDI R0 0b00000011   ; R0 = 3
LDI R1 2            ; R1 = 2 (shift amount)
SHL R0 R1 R2        ; R2 = 0b00001100 = 12
\end{lstlisting}

\textbf{Note:} Shifting left by 1 is equivalent to multiplying by 2.

\subsubsection{SHR --- Shift Right}

\begin{table}[H]
\begin{tabularx}{\textwidth}{@{}lX@{}}
\toprule
\textbf{Opcode}      & \opcode{0110} \\
\textbf{Format}      & R-Type \\
\textbf{Syntax}      & \lstinline|SHR Rs1 Rs2 Rd| \\
\textbf{Operation}   & \texttt{Rd = Rs1 >{}> Rs2} \\
\textbf{Description} & Shifts Rs1 right by Rs2 bits (logical shift). Vacated bits are filled with zeros. \\
\bottomrule
\end{tabularx}
\end{table}

\textbf{Note:} This is a logical shift (unsigned). Shifting right by 1 is equivalent to dividing by 2 (for unsigned values).

\subsection{Data Movement Instructions}

\subsubsection{MOV --- Move (Copy Register)}

\begin{table}[H]
\begin{tabularx}{\textwidth}{@{}lX@{}}
\toprule
\textbf{Opcode}      & \opcode{0111} \\
\textbf{Format}      & R-Type \\
\textbf{Syntax}      & \lstinline|MOV Rs1 Rd| \\
\textbf{Operation}   & \texttt{Rd = Rs1} \\
\textbf{Description} & Copies the value from Rs1 to Rd. Rs2 field is set to \texttt{00}. \\
\bottomrule
\end{tabularx}
\end{table}

\begin{lstlisting}
LDI R0 99      ; R0 = 99
MOV R0 R2      ; R2 = 99
\end{lstlisting}

\textbf{Encoding note:} Rs2 bits [9:8] must be \texttt{00} to select MOV (vs IN/OUT in Option~B).

\subsubsection{LDI --- Load Immediate}

\begin{table}[H]
\begin{tabularx}{\textwidth}{@{}lX@{}}
\toprule
\textbf{Opcode}      & \opcode{1010} \\
\textbf{Format}      & I-Type \\
\textbf{Syntax}      & \lstinline|LDI Rd immediate| \\
\textbf{Operation}   & \texttt{Rd = immediate} \\
\textbf{Description} & Loads an 8-bit constant value into register Rd. \\
\bottomrule
\end{tabularx}
\end{table}

\begin{lstlisting}
LDI R0 255     ; R0 = 255 (max unsigned value)
LDI R1 -1      ; R1 = 255 (same in two's complement)
LDI R2 42      ; R2 = 42
\end{lstlisting}

\subsection{Comparison Instructions}

\subsubsection{LT --- Less Than}

\begin{table}[H]
\begin{tabularx}{\textwidth}{@{}lX@{}}
\toprule
\textbf{Opcode}      & \opcode{1000} \\
\textbf{Format}      & R-Type \\
\textbf{Syntax}      & \lstinline|LT Rs1 Rs2| \\
\textbf{Operation}   & \texttt{R3 = (Rs1 < Rs2) ? 1 : 0} \\
\textbf{Description} & Compares Rs1 and Rs2 as \textbf{signed} two's complement values. Sets \reg{3} to 1 if Rs1 $<$ Rs2, otherwise 0. Rd field is ignored. \\
\bottomrule
\end{tabularx}
\end{table}

\begin{lstlisting}
LDI R0 5       ; R0 = 5
LDI R1 10      ; R1 = 10
LT R0 R1       ; R3 = 1 (because 5 < 10)
JNZ R3 loop    ; Jump taken because R3 != 0
\end{lstlisting}

\textbf{Note:} To test ``greater than'', swap the operand order: \lstinline|LT Rs2 Rs1| tests if Rs2 $<$ Rs1, which is equivalent to Rs1 $>$ Rs2.

\begin{designnote}
\textbf{Signed comparison note:} LT performs a signed comparison. Values in the range 1--100 (typical for the batch diverter project) are safe --- they are positive in both signed and unsigned interpretation. Values loaded via LDI that exceed 127 are negative in signed two's complement (e.g., \lstinline|LDI R0 200| loads the bit pattern for $-56$). This only matters if comparing mixed large values outside the normal project range.
\end{designnote}

\subsubsection{EQ --- Equal}

\begin{table}[H]
\begin{tabularx}{\textwidth}{@{}lX@{}}
\toprule
\textbf{Opcode}      & \opcode{1001} \\
\textbf{Format}      & R-Type \\
\textbf{Syntax}      & \lstinline|EQ Rs1 Rs2| \\
\textbf{Operation}   & \texttt{R3 = (Rs1 == Rs2) ? 1 : 0} \\
\textbf{Description} & Compares Rs1 and Rs2 for equality. Sets \reg{3} to 1 if equal, otherwise 0. Rd field is ignored. \\
\bottomrule
\end{tabularx}
\end{table}

\begin{lstlisting}
LDI R0 42
LDI R1 42
EQ R0 R1       ; R3 = 1 (equal)
JNZ R3 match   ; Jump taken
\end{lstlisting}

\subsection{Control Flow Instructions}

\subsubsection{JMP --- Unconditional Jump}

\begin{table}[H]
\begin{tabularx}{\textwidth}{@{}lX@{}}
\toprule
\textbf{Opcode}      & \opcode{1011} \\
\textbf{Format}      & J-Type \\
\textbf{Syntax}      & \lstinline|JMP address| \\
\textbf{Operation}   & \texttt{PC = address} \\
\textbf{Description} & Sets the program counter to the specified address (0--15). Execution continues from that instruction. Rcond field is don't-care (set to \texttt{00}). \\
\bottomrule
\end{tabularx}
\end{table}

\begin{lstlisting}
JMP 10      ; Jump to instruction at address 10
\end{lstlisting}

\subsubsection{JZ --- Jump if Zero}

\begin{table}[H]
\begin{tabularx}{\textwidth}{@{}lX@{}}
\toprule
\textbf{Opcode}      & \opcode{1100} \\
\textbf{Format}      & J-Type \\
\textbf{Syntax}      & \lstinline|JZ Rcond address| \\
\textbf{Operation}   & \texttt{if (Rcond == 0) then PC = address} \\
\textbf{Description} & Jumps to the specified address if register Rcond is zero. Otherwise, continues to the next instruction. Rcond can be any register (\reg{0}--\reg{3}). \\
\bottomrule
\end{tabularx}
\end{table}

\begin{lstlisting}
EQ R0 R1        ; R3 = 1 if equal, 0 if not
JZ R3 skip      ; Jump if NOT equal (R3 == 0)
; ... code if equal ...
; skip: ... continues here ...
\end{lstlisting}

\subsubsection{JNZ --- Jump if Not Zero}

\begin{table}[H]
\begin{tabularx}{\textwidth}{@{}lX@{}}
\toprule
\textbf{Opcode}      & \opcode{1101} \\
\textbf{Format}      & J-Type \\
\textbf{Syntax}      & \lstinline|JNZ Rcond address| \\
\textbf{Operation}   & \texttt{if (Rcond $\neq$ 0) then PC = address} \\
\textbf{Description} & Jumps to the specified address if register Rcond is not zero. Otherwise, continues to the next instruction. Rcond can be any register (\reg{0}--\reg{3}). \\
\bottomrule
\end{tabularx}
\end{table}

\begin{lstlisting}
LT R0 R1        ; R3 = 1 if R0 < R1
JNZ R3 less     ; Jump if R0 < R1
; ... code if R0 >= R1 ...
; less: ... code if R0 < R1 ...
\end{lstlisting}

\textbf{Branching on loop counters:} Because Rcond can be any register, you can branch directly on a counter register without using a comparison instruction:
\begin{lstlisting}
SUB R0 R1 R0    ; Decrement counter (R1 holds 1)
JNZ R0 loop     ; Loop while R0 != 0
\end{lstlisting}

\subsection{Miscellaneous Instructions}

\subsubsection{Reserved --- Opcode 1110}

Opcode \texttt{1110} is \textbf{reserved} for team extensions. If your team implements a custom instruction (e.g., ADDI, NOT, INC, DEC), document it in your report. If not implemented, this opcode must trigger \signal{sys\_err} as an illegal instruction (see Section~3.5).

\textbf{Extension ideas:}
\begin{itemize}[nosep]
  \item \lstinline|ADDI Rd imm| --- Add immediate to register
  \item \lstinline|NOT Rs Rd| --- Bitwise NOT
  \item \lstinline|INC Rd| / \lstinline|DEC Rd| --- Increment/decrement
  \item \lstinline|SWAPNIB Rd| --- Swap nibbles (rotate by 4)
\end{itemize}

\textbf{Note:} If you need a no-operation, use \lstinline|MOV R0 R0| (copies \reg{0} to itself).

\subsubsection{HALT --- Halt Execution}

\begin{table}[H]
\begin{tabularx}{\textwidth}{@{}lX@{}}
\toprule
\textbf{Opcode}      & \opcode{1111} \\
\textbf{Format}      & --- \\
\textbf{Syntax}      & \lstinline|HALT| \\
\textbf{Operation}   & Sets \signal{cpu\_halt} latch \\
\textbf{Description} & Sets the internal \signal{cpu\_halt} flag, which stops the program counter from advancing. No further instructions execute. \\
\bottomrule
\end{tabularx}
\end{table}

\textbf{Encoding:} \texttt{1111\;0000\;0000\;0000} (all zeros after opcode)

\textbf{Platform behaviour:}
\begin{itemize}[nosep]
  \item HALT sets \signal{cpu\_halt} $\rightarrow$ PC frozen, registers and output port preserved
  \item Platform treats \signal{cpu\_halt} as a halt source: $\text{HALTED} = \text{STOP} \mid \text{FAULT} \mid \text{SYS\_ERR} \mid \text{CPU\_HALT}$
  \item \signal{motor\_en} is forced to 0 (motor safety interlock, see Section~\ref{sec:haltlogic})
  \item \signal{alarm} is \textbf{not} asserted (HALT is intentional, not a fault)
  \item \textbf{Recovery:} Only \signal{reboot} clears \signal{cpu\_halt} and restarts execution from address 0
\end{itemize}

\textbf{Note:} In normal batch diverter operation, the program loops infinitely and HALT is never reached. HALT is useful for test programs and as a safety backstop.

\subsection{I/O Instructions (Option B)}\label{sec:io_instructions}

If your team selects \textbf{Option B: Dedicated I/O Instructions}, implement the following. If using Option~A or~C, skip this section.

\subsubsection{IN --- Read Input Port}

\begin{table}[H]
\begin{tabularx}{\textwidth}{@{}lX@{}}
\toprule
\textbf{Opcode}      & Shares with MOV: \opcode{0111} with Rs2 = \texttt{11} \\
\textbf{Format}      & R-Type (special) \\
\textbf{Syntax}      & \lstinline|IN Rd| \\
\textbf{Operation}   & \texttt{Rd = input\_port} (clears \signal{item\_pulse} latch) \\
\textbf{Description} & Reads the current state of the input port into register Rd. The \signal{item\_pulse} latch is cleared on read (see Section~\ref{sec:itempulse}). \\
\bottomrule
\end{tabularx}
\end{table}

\textbf{Encoding:}

\begin{center}
\begin{bytefield}[bitwidth=1.6em]{16}
  \bitheader{0-15} \\
  \bitbox{4}{\footnotesize\texttt{0111}} &
  \bitbox{2}{\footnotesize \texttt{00}} &
  \bitbox{2}{\footnotesize \texttt{11}} &
  \bitbox{2}{\footnotesize Rd} &
  \bitbox{6}{\footnotesize \texttt{000000}}
\end{bytefield}
\end{center}

The \texttt{Rs2 = 11} distinguishes IN from MOV (where Rs2 = 00).

\begin{lstlisting}
IN R0          ; R0 = input port state (bit 0 = ITEM_PULSE)
LDI R1 0x01   ; Mask for ITEM_PULSE (bit 0)
AND R0 R1 R0  ; R0 = 1 if item detected, else 0
\end{lstlisting}

\subsubsection{OUT --- Write Output Port}

\begin{table}[H]
\begin{tabularx}{\textwidth}{@{}lX@{}}
\toprule
\textbf{Opcode}      & Shares with MOV: \opcode{0111} with Rs2 = \texttt{10} \\
\textbf{Format}      & R-Type (special) \\
\textbf{Syntax}      & \lstinline|OUT Rs| \\
\textbf{Operation}   & \texttt{output\_port = Rs} \\
\textbf{Description} & Writes the contents of Rs to the output port (\signal{gate\_sel}, \signal{motor\_en}, \signal{status}). \\
\bottomrule
\end{tabularx}
\end{table}

\textbf{Encoding:}

\begin{center}
\begin{bytefield}[bitwidth=1.6em]{16}
  \bitheader{0-15} \\
  \bitbox{4}{\footnotesize\texttt{0111}} &
  \bitbox{2}{\footnotesize Rs} &
  \bitbox{2}{\footnotesize \texttt{10}} &
  \bitbox{2}{\footnotesize \texttt{00}} &
  \bitbox{6}{\footnotesize \texttt{000000}}
\end{bytefield}
\end{center}

The \texttt{Rs2 = 10} distinguishes OUT from MOV and IN.

\begin{lstlisting}
LDI R2 0x02   ; MOTOR_EN = 1, GATE_SEL = 0 (Gate A)
OUT R2         ; Write to output port
\end{lstlisting}

\subsubsection{I/O Encoding Summary (Option B)}

\begin{table}[H]
\centering
\begin{tabular}{@{}clll@{}}
\toprule
\textbf{Rs2} & \textbf{Instruction} & \textbf{Operation} \\
\midrule
\texttt{00} & \lstinline|MOV Rs1 Rd| & Rd = Rs1 \\
\texttt{01} & (undefined --- triggers \signal{sys\_err}) & --- \\
\texttt{10} & \lstinline|OUT Rs1| & output\_port = Rs1 \\
\texttt{11} & \lstinline|IN Rd| & Rd = input\_port \\
\bottomrule
\end{tabular}
\end{table}

\begin{designnote}
By overloading the MOV opcode with Rs2 variants, we avoid consuming additional opcodes while providing clean I/O semantics. The decoder checks Rs2 to select between register move and port operations. The \texttt{Rs2=01} encoding is undefined and must trigger \signal{sys\_err} if executed (see Section~3.5).
\end{designnote}

% ══════════════════════════════════════════════════════════════════════
\section{Instruction Encoding Summary}

\begin{table}[H]
\centering
\small
\begin{tabular}{@{}cclll@{}}
\toprule
\textbf{Opcode} & \textbf{Binary} & \textbf{Mnemonic} & \textbf{Format} & \textbf{Operation} \\
\midrule
0  & \texttt{0000} & ADD  & R & Rd = Rs1 + Rs2 \\
1  & \texttt{0001} & SUB  & R & Rd = Rs1 $-$ Rs2 \\
2  & \texttt{0010} & AND  & R & Rd = Rs1 \& Rs2 \\
3  & \texttt{0011} & OR   & R & Rd = Rs1 $|$ Rs2 \\
4  & \texttt{0100} & XOR  & R & Rd = Rs1 $\oplus$ Rs2 \\
5  & \texttt{0101} & SHL  & R & Rd = Rs1 $\ll$ Rs2 \\
6  & \texttt{0110} & SHR  & R & Rd = Rs1 $\gg$ Rs2 \\
7  & \texttt{0111} & MOV  & R & Rd = Rs1 \quad (Rs2=00) \\
7  & \texttt{0111} & OUT  & R & output = Rs1 \quad (Rs2=10)\textsuperscript{\dag} \\
7  & \texttt{0111} & IN   & R & Rd = input \quad (Rs2=11)\textsuperscript{\dag} \\
8  & \texttt{1000} & LT   & R & R3 = (Rs1 $<$ Rs2) \\
9  & \texttt{1001} & EQ   & R & R3 = (Rs1 $=$ Rs2) \\
10 & \texttt{1010} & LDI  & I & Rd = immediate \\
11 & \texttt{1011} & JMP  & J & PC = address \\
12 & \texttt{1100} & JZ   & J & if Rcond$=$0: PC = addr \\
13 & \texttt{1101} & JNZ  & J & if Rcond$\neq$0: PC = addr \\
14 & \texttt{1110} & ---  & --- & Reserved (extension or \signal{sys\_err}) \\
15 & \texttt{1111} & HALT & --- & Set \signal{cpu\_halt} \\
\bottomrule
\end{tabular}
\caption*{\textsuperscript{\dag} Option B only. See Section~\ref{sec:io_instructions} for I/O instruction details.}
\end{table}

% ══════════════════════════════════════════════════════════════════════
\section{Binary Encoding Examples}

\subsection{Example 1: ADD R0 R1 R2}

\begin{lstlisting}[language={}]
ADD  R0   R1   R2   unused
0000  00   01   10   000000

Binary: 0000 0001 1000 0000
Hex:    0x0180
\end{lstlisting}

\subsection{Example 2: LDI R0 42}

\begin{lstlisting}[language={}]
LDI  R0   --   immediate(42)
1010  00   00   00101010

Binary: 1010 0000 0010 1010
Hex:    0xA02A
\end{lstlisting}

\subsection{Example 3: JNZ R3 5}

\begin{lstlisting}[language={}]
JNZ  Rcond(R3) unused   address(5)
1101    11     000000    0101

Binary: 1101 1100 0000 0101
Hex:    0xDC05
\end{lstlisting}

\subsection{Example 4: JZ R0 3}

\begin{lstlisting}[language={}]
JZ   Rcond(R0) unused   address(3)
1100    00     000000    0011

Binary: 1100 0000 0000 0011
Hex:    0xC003
\end{lstlisting}

\subsection{Example 5: JMP 10}

\begin{lstlisting}[language={}]
JMP  Rcond(--) unused   address(10)
1011    00     000000    1010

Binary: 1011 0000 0000 1010
Hex:    0xB00A
\end{lstlisting}

\subsection{Example 6: LT R0 R1 (compare)}

\begin{lstlisting}[language={}]
LT   R0   R1   --   unused
1000  00   01   00   000000

Binary: 1000 0001 0000 0000
Hex:    0x8100
\end{lstlisting}

\subsection{Example 7: IN R0 (Option B)}

\begin{lstlisting}[language={}]
IN   --   11   R0   unused
0111  00   11   00   000000

Binary: 0111 0011 0000 0000
Hex:    0x7300
\end{lstlisting}

\subsection{Example 8: OUT R2 (Option B)}

\begin{lstlisting}[language={}]
OUT  R2   10   --   unused
0111  10   10   00   000000

Binary: 0111 1010 0000 0000
Hex:    0x7A00
\end{lstlisting}

% ══════════════════════════════════════════════════════════════════════
\section{Program Memory}

The EE8 uses a ROM (Read-Only Memory) to store programs.

\begin{table}[H]
\centering
\begin{tabular}{@{}ll@{}}
\toprule
\textbf{Parameter} & \textbf{Value} \\
\midrule
Word size      & 16 bits (one instruction) \\
Address width  & 4 bits \\
Capacity       & 16 words \\
Address range  & 0--15 (\texttt{0x0}--\texttt{0xF}) \\
\bottomrule
\end{tabular}
\end{table}

\subsection{Program Counter (PC)}

\begin{itemize}[nosep]
  \item 4-bit register
  \item Initialises to 0 on reset
  \item Increments by 1 after each instruction (unless jump taken)
  \item Wraps around from 15 to 0
  \item Frozen when \signal{halted} (see Section~\ref{sec:haltlogic})
\end{itemize}

\subsection{Loading Programs}

Students will build two ROM modules:
\begin{enumerate}[nosep]
  \item \textbf{External ROM} --- Instructor-provided test programs (pre-loaded)
  \item \textbf{Internal ROM} --- Student-built program memory
\end{enumerate}

A \textbf{LOAD signal} transfers the test program from external ROM to internal ROM before execution begins.

\subsection{Fetch-Decode-Execute Cycle}

\begin{enumerate}[nosep]
  \item \textbf{Fetch:} Read instruction from ROM at address PC
  \item \textbf{Decode:} Extract opcode and operands
  \item \textbf{Execute:} Perform operation
  \item \textbf{Update PC:} PC = PC + 1 (or jump target if branch taken)
\end{enumerate}

% ══════════════════════════════════════════════════════════════════════
\section{Timing and Clocking}

The CPU operates on a single clock signal in a \textbf{single-cycle architecture}.

\begin{table}[H]
\centering
\begin{tabularx}{\textwidth}{@{}lX@{}}
\toprule
\textbf{Mode} & \textbf{Description} \\
\midrule
Step mode & Manual clock button for debugging --- one instruction per press \\
Run mode  & Continuous clock signal from oscillator \\
\bottomrule
\end{tabularx}
\end{table}

\subsection{Clock Phases (single-cycle)}

On each rising clock edge:
\begin{enumerate}[nosep]
  \item Current instruction executes (fetch, decode, execute happen combinationally)
  \item Results written to registers
  \item PC updates
\end{enumerate}

All operations are \textbf{atomic} --- each instruction completes in a single clock cycle.

\begin{designnote}
Single-cycle architecture is simpler to design and understand than pipelined CPUs, making it ideal for educational purposes. The trade-off is that all instructions take the same time, even if some could be faster.
\end{designnote}

% ══════════════════════════════════════════════════════════════════════
\section{Reset State Summary}

On \signal{reboot} (or initial power-on), the CPU enters the following known state:

\begin{table}[H]
\centering
\begin{tabular}{@{}lll@{}}
\toprule
\textbf{Element} & \textbf{Reset Value} & \textbf{Notes} \\
\midrule
\reg{0}          & \texttt{0x00} & \\
\reg{1}          & \texttt{0x00} & \\
\reg{2}          & \texttt{0x00} & \\
\reg{3}          & \texttt{0x00} & \\
PC               & \texttt{0x0}  & Execution starts at address 0 \\
Output port      & \texttt{0x00} & \signal{motor\_en}=0, \signal{gate\_sel}=0 (Gate A) \\
\signal{sys\_err}  & \texttt{0}    & Cleared \\
\signal{cpu\_halt} & \texttt{0}    & Cleared \\
\bottomrule
\end{tabular}
\end{table}

% ══════════════════════════════════════════════════════════════════════
\section{Assembly Language Syntax}

\subsection{Comments}

\begin{lstlisting}
; This is a comment
ADD R0 R1 R2    ; Inline comment
\end{lstlisting}

\subsection{Labels}

\begin{lstlisting}
start:  LDI R0 10
        LDI R1 1
loop:   SUB R0 R1 R0
        JNZ R0 loop    ; Jump to 'loop' if R0 != 0
        HALT
\end{lstlisting}

\subsection{Numeric Formats}

\begin{lstlisting}
LDI R0 42          ; Decimal
LDI R0 0x2A        ; Hexadecimal
LDI R0 0b00101010  ; Binary
\end{lstlisting}

% ══════════════════════════════════════════════════════════════════════
\section{Example Programs}

\subsection{Batch Diverter: $A \neq B$ with Full Alternation (Option B)}

This is the primary example: a correct batch diverter that sends $A=5$ items to Gate~A, then $B=3$ items to Gate~B, repeating forever. Uses all 16 instruction slots.

\begin{lstlisting}
; Batch diverter: A=5 to Gate A, B=3 to Gate B, alternating forever
; Option B I/O. All 16 instruction slots used.
;
; Register usage:
;   R0 = item counter (counts down to 0)
;   R1 = scratch (input, constants)
;   R2 = output port value (bit 0 = GATE_SEL, bit 1 = MOTOR_EN)
;   R3 = scratch (masks)
;
; After REBOOT: all regs = 0, PC = 0.

; === Setup (runs once at start, also re-executed on Gate A reload) ===
0:  LDI R0 5            ; R0 = batch count A (5 items)
1:  LDI R2 0x02         ; R2 = MOTOR_EN=1, GATE_SEL=0 (Gate A)

; === Counting loop (shared for both gates) ===
2:  OUT R2              ; Drive output: motor on, current gate
3:  IN R1               ; Read input port (clears ITEM_PULSE latch)
4:  LDI R3 0x01         ; ITEM_PULSE mask (bit 0)
5:  AND R1 R3 R1        ; Isolate ITEM_PULSE
6:  JZ R1 3             ; No item detected? Keep polling
7:  LDI R1 1            ; Decrement constant
8:  SUB R0 R1 R0        ; R0 = R0 - 1
9:  JNZ R0 3            ; More items in batch? Continue counting

; === Batch complete: toggle gate and reload count ===
10: LDI R1 0x01         ; GATE_SEL toggle mask (bit 0)
11: XOR R2 R1 R2        ; Flip GATE_SEL: Gate A (0x02) <-> Gate B (0x03)
12: AND R2 R1 R3        ; R3 = new GATE_SEL (0 = Gate A, 1 = Gate B)
13: LDI R0 3            ; Load count B (optimistic for Gate B)
14: JNZ R3 2            ; If Gate B selected (R3=1) -> counting loop
15: LDI R0 5            ; Gate A selected: load count A = 5
                         ; PC wraps 15 -> 0
\end{lstlisting}

\textbf{Execution trace (first 3 batches):}

\begin{table}[H]
\centering
\small
\begin{tabular}{@{}llcccc@{}}
\toprule
\textbf{Phase} & \textbf{Addresses} & \textbf{R0} & \textbf{R2} & \textbf{Gate} & \textbf{Items} \\
\midrule
Setup   & $0 \rightarrow 1$                   & 5         & \texttt{0x02} & A & --- \\
Batch 1 & $2 \rightarrow \cdots \rightarrow 9$ & $5 \rightarrow 0$ & \texttt{0x02} & A & 5 \\
Toggle  & $10 \rightarrow \cdots \rightarrow 14 \rightarrow 2$ & 3 & \texttt{0x03} & B & --- \\
Batch 2 & $2 \rightarrow \cdots \rightarrow 9$ & $3 \rightarrow 0$ & \texttt{0x03} & B & 3 \\
Toggle  & $10 \rightarrow \cdots \rightarrow 15 \rightarrow 0 \rightarrow 1 \rightarrow 2$ & 5 & \texttt{0x02} & A & --- \\
Batch 3 & $2 \rightarrow \cdots \rightarrow 9$ & $5 \rightarrow 0$ & \texttt{0x02} & A & 5 \\
\ldots  & repeats & & & & \\
\bottomrule
\end{tabular}
\end{table}

\textbf{Pattern:} 5(A), 3(B), 5(A), 3(B), \ldots\ $\checkmark$

\subsection{Simple Batch Diverter: Equal Counts (Option B)}

A simpler version where both gates get the same count. Useful as a starting point.

\begin{lstlisting}
; Simple batch diverter: 5 items per gate, alternating
; Option B I/O. 12 instructions.

; === Setup ===
0:  LDI R2 0x02         ; MOTOR_EN=1, GATE_SEL=0 (Gate A)
1:  LDI R0 5            ; Batch count = 5

; === Counting loop ===
2:  OUT R2              ; Drive outputs
3:  IN R1               ; Read input
4:  LDI R3 0x01         ; ITEM_PULSE mask
5:  AND R1 R3 R1        ; Isolate pulse
6:  JZ R1 3             ; Poll until item arrives
7:  LDI R1 1            ; Decrement constant
8:  SUB R0 R1 R0        ; R0--
9:  JNZ R0 3            ; Count more items

; === Toggle and restart ===
10: LDI R1 0x01         ; Toggle mask
11: XOR R2 R1 R2        ; Flip GATE_SEL
                         ; PC falls through to 12
12: LDI R0 5            ; Reload same count
13: JMP 2               ; Back to counting loop

14: HALT                ; Never reached
15: HALT
\end{lstlisting}

\subsection{Arithmetic Example: Add Two Numbers}

\begin{lstlisting}
; Adds 25 + 17, stores result in R2
; 4 instructions
0:  LDI R0 25       ; R0 = 25
1:  LDI R1 17       ; R1 = 17
2:  ADD R0 R1 R2    ; R2 = 42
3:  HALT
\end{lstlisting}

\subsection{Find Maximum of Two Numbers}

\begin{lstlisting}
; Finds max of R0 and R1, stores in R2
; 8 instructions
0:  LDI R0 45           ; R0 = 45
1:  LDI R1 72           ; R1 = 72
2:  LT R0 R1            ; R3 = 1 if R0 < R1
3:  JNZ R3 6            ; If R0 < R1, jump to r1_bigger
4:  MOV R0 R2           ; R0 is bigger or equal
5:  JMP 7               ; Skip to done
6:  MOV R1 R2           ; R1 is bigger
7:  HALT                ; R2 = 72 (the maximum)
\end{lstlisting}

\subsection{Polling Input Example (Option B)}

\begin{lstlisting}
; Wait for ITEM_PULSE, then turn off motor
; Demonstrates I/O polling pattern
; 7 instructions

0:  LDI R2 0x02         ; MOTOR_EN=1, GATE_SEL=0
1:  OUT R2              ; Enable motor
2:  IN R0               ; Read inputs
3:  LDI R1 0x01         ; ITEM_PULSE mask
4:  AND R0 R1 R0        ; Isolate bit 0
5:  JZ R0 2             ; No pulse? Keep polling
6:  HALT                ; Item detected, stop
\end{lstlisting}

\subsection{Countdown Example}

\begin{lstlisting}
; Counts down from 10 to 0 using the JNZ Rcond syntax
; Demonstrates branching on a counter register directly
; 4 instructions

0:  LDI R0 10           ; R0 = counter = 10
1:  LDI R1 1            ; R1 = decrement constant
2:  SUB R0 R1 R0        ; R0 = R0 - 1
3:  JNZ R0 2            ; Loop while R0 != 0
4:  HALT                ; R0 = 0, done
\end{lstlisting}

% ══════════════════════════════════════════════════════════════════════
\appendix
\section{Quick Reference Card}

\begin{tcolorbox}[colback=codebg, colframe=ee8blue, title={\large\bfseries EE8 CPU Quick Reference --- Version 4.1}, fonttitle=\color{white}]
\small
\begin{multicols}{2}

\textbf{Registers}\\[2pt]
\begin{tabular}{@{}ll@{}}
\reg{0}, \reg{1} & General purpose \\
\reg{2} & General purpose \\
\reg{3} & General purpose + comparison dest \\
\end{tabular}

\vspace{6pt}
\textbf{Arithmetic}\\[2pt]
\begin{tabular}{@{}l@{}}
\lstinline|ADD Rs1 Rs2 Rd| \\
\lstinline|SUB Rs1 Rs2 Rd| \\
\end{tabular}

\vspace{6pt}
\textbf{Bitwise}\\[2pt]
\begin{tabular}{@{}l@{}}
\lstinline|AND Rs1 Rs2 Rd| \\
\lstinline|OR  Rs1 Rs2 Rd| \\
\lstinline|XOR Rs1 Rs2 Rd| \\
\end{tabular}

\vspace{6pt}
\textbf{Shift}\\[2pt]
\begin{tabular}{@{}l@{}}
\lstinline|SHL Rs1 Rs2 Rd| \\
\lstinline|SHR Rs1 Rs2 Rd| \\
\end{tabular}

\vspace{6pt}
\textbf{Compare} (result $\rightarrow$ \reg{3})\\[2pt]
\begin{tabular}{@{}l@{}}
\lstinline|LT Rs1 Rs2| \quad (R3 = Rs1$<$Rs2) \\
\lstinline|EQ Rs1 Rs2| \quad (R3 = Rs1$=$Rs2) \\
\end{tabular}

\columnbreak

\textbf{Data Movement}\\[2pt]
\begin{tabular}{@{}ll@{}}
\lstinline|MOV Rs Rd|      & Copy register \\
\lstinline|LDI Rd imm|     & Load immediate (0--255) \\
\lstinline|IN Rd|           & Read input port (Opt.\ B) \\
\lstinline|OUT Rs|          & Write output port (Opt.\ B) \\
\end{tabular}

\vspace{6pt}
\textbf{Control Flow} (addr: 0--15)\\[2pt]
\begin{tabular}{@{}ll@{}}
\lstinline|JMP addr|        & Unconditional jump \\
\lstinline|JZ  Rcond addr|  & Jump if Rcond $= 0$ \\
\lstinline|JNZ Rcond addr|  & Jump if Rcond $\neq 0$ \\
\end{tabular}

\vspace{6pt}
\textbf{Miscellaneous}\\[2pt]
\begin{tabular}{@{}ll@{}}
\lstinline|HALT|            & Stop execution \\
(opcode \texttt{1110})      & Reserved / \signal{sys\_err} \\
\end{tabular}

\vspace{6pt}
\textbf{I/O Port Bits}\\[2pt]
\begin{tabular}{@{}ll@{}}
Input [0]  & \signal{item\_pulse} (cleared on read) \\
Output [0] & \signal{gate\_sel} \\
Output [1] & \signal{motor\_en} \\
Output [2+]& \signal{status} (optional) \\
\end{tabular}

\end{multicols}

\vspace{4pt}
\textbf{Platform Signals} (hardware, NOT software-polled)\\[2pt]
\begin{tabular}{@{}ll@{}}
\signal{reboot}   & Async reset (clears all state, PC=0) \\
\signal{stop}     & Pause (clock gated, state preserved) \\
\signal{fault}    & Pause + \signal{alarm} \\
\signal{sys\_err}  & Illegal insn $\rightarrow$ halt + \signal{alarm} (latched) \\
\signal{cpu\_halt} & HALT insn $\rightarrow$ halt, no \signal{alarm} (latched) \\
\end{tabular}

\vspace{4pt}
$\text{HALTED} = \text{STOP} \mid \text{FAULT} \mid \text{SYS\_ERR} \mid \text{CPU\_HALT}$\\
$\text{MOTOR\_EN}_\text{pin} = \text{MOTOR\_EN}_\text{cpu} \wedge \overline{\text{HALTED}}$\\
$\text{ALARM} = \text{FAULT} \mid \text{SYS\_ERR}$\\[4pt]
Memory: 16 instructions (addresses 0--F). PC wraps from 15 $\rightarrow$ 0.
\end{tcolorbox}

% ══════════════════════════════════════════════════════════════════════
\section{Opcode Map}

\begin{table}[H]
\centering
\begin{tabular}{@{}cl|cl@{}}
\toprule
\textbf{Binary} & \textbf{Mnemonic} & \textbf{Binary} & \textbf{Mnemonic} \\
\midrule
\texttt{0000} & ADD         & \texttt{1000} & LT \\
\texttt{0001} & SUB         & \texttt{1001} & EQ \\
\texttt{0010} & AND         & \texttt{1010} & LDI \\
\texttt{0011} & OR          & \texttt{1011} & JMP \\
\texttt{0100} & XOR         & \texttt{1100} & JZ \\
\texttt{0101} & SHL         & \texttt{1101} & JNZ \\
\texttt{0110} & SHR         & \texttt{1110} & (reserved / \signal{sys\_err}) \\
\texttt{0111} & MOV/IN/OUT  & \texttt{1111} & HALT \\
\bottomrule
\end{tabular}
\end{table}

\textbf{MOV/IN/OUT disambiguation (opcode \texttt{0111}):}

\begin{table}[H]
\centering
\begin{tabular}{@{}cl@{}}
\toprule
\textbf{Rs2} & \textbf{Instruction} \\
\midrule
\texttt{00} & \lstinline|MOV Rs1 Rd| \\
\texttt{01} & (undefined $\rightarrow$ \signal{sys\_err}) \\
\texttt{10} & \lstinline|OUT Rs1| \\
\texttt{11} & \lstinline|IN Rd| \\
\bottomrule
\end{tabular}
\end{table}

\textbf{J-Type format (JMP/JZ/JNZ):}\\
\texttt{[15:12] opcode \quad [11:10] Rcond \quad [9:4] unused \quad [3:0] address}
\begin{itemize}[nosep]
  \item JMP: Rcond ignored (set to \texttt{00})
  \item JZ: jump if Rcond $= 0$
  \item JNZ: jump if Rcond $\neq 0$
\end{itemize}

% ══════════════════════════════════════════════════════════════════════
\section{I/O Interface Options Summary}

\begin{table}[H]
\centering
\begin{tabularx}{\textwidth}{@{}clXX@{}}
\toprule
\textbf{Option} & \textbf{Mechanism} & \textbf{Pros} & \textbf{Cons} \\
\midrule
A & Memory-mapped   & No new instructions          & Loses ROM addresses \\
B & IN/OUT instructions & Clean, explicit          & Uses opcode space \\
C & Register-mapped & Simplest hardware            & Loses registers \\
\bottomrule
\end{tabularx}
\end{table}


\textbf{I/O Port Signals (all options):}

\begin{table}[H]
\centering
\begin{tabular}{@{}cclp{6cm}@{}}
\toprule
\textbf{Port} & \textbf{Bit} & \textbf{Signal} & \textbf{Notes} \\
\midrule
Input  & 0  & \signal{item\_pulse} & Cleared on read \\
Output & 0  & \signal{gate\_sel}   & 0 = Gate A, 1 = Gate B \\
Output & 1  & \signal{motor\_en}   & Request (overridden when \signal{halted}) \\
Output & 2+ & \signal{status}      & Optional debug bits \\
\bottomrule
\end{tabular}
\end{table}

\textbf{Not in I/O ports} (hardware-driven): \signal{alarm}, \signal{stop}, \signal{fault}, \signal{reboot}

\section{Changes from Previous Versions}

\subsection*{Version 4.1 (correction)}
\begin{itemize}[nosep]
  \item \textbf{I-Type encoding fixed}: Rd is at bits [11:10], not [9:8]. Bytefield diagram, field list, and Example~2 corrected to match reference CPU implementation and Count-to-5 handout.
\end{itemize}

\subsection*{Version 4.0 (platform signals, Rcond branches, safety)}
\begin{itemize}[nosep]
  \item \textbf{New Section 3: Platform Signals \& Safety Logic} --- REBOOT, STOP, FAULT, SYS\_ERR, CPU\_HALT fully specified as hardware signals, not software-polled I/O
  \item \textbf{J-Type format changed}: added Rcond field [11:10] --- JZ/JNZ can now test any register, not just \reg{3}
  \item \textbf{HALT behaviour clarified}: sets \signal{cpu\_halt} latch, \signal{motor\_en} forced 0, \signal{alarm} not asserted, only cleared by \signal{reboot}
  \item \textbf{SYS\_ERR fully specified}: triggers on undefined opcode or undefined Rs2=01 sub-encoding, latched until \signal{reboot}
  \item \textbf{Platform halt logic}: $\text{HALTED} = \text{STOP} \mid \text{FAULT} \mid \text{SYS\_ERR} \mid \text{CPU\_HALT}$
  \item \textbf{Alarm logic}: $\text{ALARM} = \text{FAULT} \mid \text{SYS\_ERR}$ (hardware-driven, removed from CPU output port)
  \item \textbf{I/O port revised}: input carries only \signal{item\_pulse} (bit 0), STOP/FAULT/REBOOT removed from input port
  \item \textbf{ITEM\_PULSE semantics}: hardware-latched, cleared on read (prevents multi-counting)
  \item \textbf{Reset state specified}: Section~11 documents all state after \signal{reboot}
  \item \textbf{Signed comparison note} added for LT instruction
  \item \textbf{Section numbering fixed} throughout
  \item \textbf{All examples corrected}: fixed register name typos, fixed polling/counting loop bugs, all examples use new JZ/JNZ Rcond syntax
  \item \textbf{New example 13.1}: correct $A \neq B$ batch diverter with full execution trace
\end{itemize}

\subsection*{Version 3.0 (I/O and batch diverter)}
\begin{itemize}[nosep]
  \item Added I/O interface section (Options A, B, C)
  \item Added IN and OUT instructions (Option B) via MOV opcode overloading
  \item Renamed RO $\rightarrow$ R2, RF $\rightarrow$ R3 for uniformity (all registers now general-purpose)
  \item Removed NOP (opcode 1110 now reserved for team extensions)
  \item Updated examples for batch diverter scenario
  \item Added I/O port bit assignments for station signals
  \item Aligned with EG2300 project specification
\end{itemize}

\subsection*{Version 2.0 (4-bit addressing)}
\begin{itemize}[nosep]
  \item Reduced program memory from 256 to 16 instructions
  \item Changed PC from 8-bit to 4-bit counter
  \item Updated J-type format: address field now [3:0] (4 bits)
  \item All jump/branch addresses now 0--15 range
  \item Added design rationale for comparison register approach
\end{itemize}

\vfill
\begin{center}
\textcolor{ee8gray}{\rule{0.4\textwidth}{0.4pt}}\\[6pt]
\textit{End of EE8 CPU ISA Specification}
\end{center}

\end{document}
